\sscp{\tsc{North Halmahera}}{C-10-02}
The following are all Papuan languages of the North Halmahera family.
See \srf{14-03-02} for more discussion of these terms
and the justification for reconstructing *kusor.
\\

\noindent{\wg\st{2pt}\begin{tabularx}{\linewidth}
{lllll>{\raggedright\arraybackslash}X}
\lsptoprule
%Language	&	ISO	&	Term	&	gloss	&	Etymon	&	Source	\\	\midrule
%\endfirsthead
Language	&	ISO	&	Term	&	gloss	&	Etymon	&	Source	\\	\midrule
%\endhead
\ili{Sahu}		&	saj	&	ʔúsoro	&	cuscus	&	*kusor	&	\cite[131]{VisserVoorhoeve1987}	\\
\ili{Ternate}	&	tft	&	kuso	&	cuscus	&	*kusor	&	\cite[269]{Schapper2011}	\\
\ili{Tidore}	&	tvo	&	kuso	&	cuscus	&	*kusor	&	\cite[269]{Schapper2011}	\\
\ili{Galela}	&	gbi	&	kuso	&	cuscus	&	*kusor	&	\cite[223]{vanBaarda1895}	\\
\ili{Tobelo}	&	tlb	&	kuho	&	cuscus	&	*kusor	&	\cite[212]{Hueting1908}	\\
\ili{Pagu}	&	pgu	&	kuso	&	cuscus	&	*kusor	&	Dalan Mehuli Perangin Angin pers. comm. May 2023	\\
W. \ili{Makian}	&	mky	&	kusu	&	cuscus	&	*kusor	&	\cite[81]{Collins1982b}	\\
\lspbottomrule
\end{tabularx}}

%\noindent{\wg\st{2pt}\begin{xltabular}{\linewidth}
%{lllll>{\raggedright\arraybackslash}X}
%\lsptoprule
%
%\lspbottomrule
%\end{xltabular}}